\documentclass[conference,onecolumn]{IEEEtran}

\usepackage[utf8]{inputenc}
\usepackage[T1]{fontenc}
\usepackage{booktabs}
\usepackage{hyperref}
\usepackage{orcidlink}

\hypersetup{colorlinks=false, pdfborder={0 0 0}}

\begin{document}

\title{Geometry-resolved Characterization of a Dual-detector MA-XRF Scanner}

\author{
\small
\begin{tabular}{@{}c@{\hspace{1.2em}}c@{\hspace{1.2em}}c@{}}
\begin{tabular}[t]{@{}c@{}}
\textbf{Dimitrije Pešić}\\
\textit{University of Belgrade}\\
\textit{School of Electrical Engineering}\\
Belgrade, Serbia\\
pd230014d@student.etf.bg.ac.rs
\end{tabular}
&
\begin{tabular}[t]{@{}c@{}}
\textbf{Janko Vukobratović}\\
\textit{University of Belgrade}\\
\textit{School of Electrical Engineering}\\
Belgrade, Serbia\\
vj240035d@student.etf.bg.ac.rs
\end{tabular}
&
\begin{tabular}[t]{@{}c@{}}
\textbf{Aleksandra Stojanović}\\
\textit{University of Belgrade}\\
\textit{School of Electrical Engineering}\\
Belgrade, Serbia\\
sa230234d@student.etf.bg.ac.rs
\end{tabular}
\\
\addlinespace[0.6em]
\begin{tabular}[t]{@{}c@{}}
\textbf{Giulia Ristori}\\
\textit{Ars Mensurae}\\
Rome, Italy\\
giulia.ristori@arsmensurae.it
\end{tabular}
&
\begin{tabular}[t]{@{}c@{}}
\textbf{Stefano Ridolfi}\\
\textit{IDArtScience Srl}\\
Rome, Italy\\
stefano@arsmensurae.it\\
\orcidlink{0000-0002-2866-6344} ORCID: 0000-0002-2866-6344
\end{tabular}
&
\begin{tabular}[t]{@{}c@{}}
\textbf{Maja Gajić-Kvaščev}\\
\textit{VINARH Center}\\
\textit{University of Belgrade}\\
\textit{Vinča Institute of Nuclear Sciences}\\
Belgrade, Serbia\\
gajicm@vin.bg.ac.rs\\
\orcidlink{0000-0003-1288-4248} ORCID: 0000-0003-1288-4248
\end{tabular}
\\
\addlinespace[0.6em]
\multicolumn{3}{c}{
\begin{tabular}[t]{@{}c@{}}
\textbf{Goran Kvaščev}\\
\textit{Signals and systems Department}\\
\textit{University of Belgrade}\\
\textit{School of Electrical Engineering}\\
Belgrade, Serbia\\
kvascev@etf.bg.ac.rs\\
\orcidlink{0000-0001-8642-0361} ORCID: 0000-0001-8642-0361
\end{tabular}
}
\end{tabular}
}

\maketitle

\bigskip

\noindent
Macro-X-ray fluorescence (MA-XRF) scanners with two detectors in the setup
are used so that their signals are normally summed to improve the
signal-to-noise ratio, and the difference between them has been neglected.
We show that this difference, plus one extra measurement, scanning the same
painting again with the canvas tilted forward, is enough to fully
characterize the instrument, with no dedicated calibration measurements.
Tilting the canvas changes the path the fluorescence photons travel to each
detector but leaves the detectors themselves unchanged. Comparing the two
scans therefore splits the observed differences into two parts: detector
properties, which do not depend on tilt, and geometric effects, which change
with tilt in a predictable, energy-dependent way. From two routine scans of
one painting, we obtain per-element efficiency ratios with uncertainties,
the effective viewing angle of each detector, and a flat-field correction
map, and we measure how sensitive the element maps are to canvas
positioning. The result is an energy-weighted fusion of the two detectors
that outperforms simple summing, and a practical estimate of the error
caused by imperfect mounting.

\end{document}
